\section{Fichamento 1: Ahmed e Wahed (2020)}

\noindent\textbf{Autoria do fichamento:} Eduardo Velloso Landeiro

\subsection*{1. Referência bibliográfica}
\noindent\fullcite{ahmed2020}

\subsection*{2. Tema do artigo}
A concentração da pesquisa em inteligência artificial (IA) em poucas organizações depois que o
aprendizado profundo ganhou força, e o papel do acesso desigual a poder computacional
(\emph{compute}) nessa concentração. O artigo compara grandes empresas de tecnologia,
universidades de elite e universidades de menor prestígio pela presença nas principais
conferências da área.

\subsection*{3. Problema investigado}
A suspeita de que a IA moderna, por depender de \emph{hardware} caro, virou um campo de quem tem e
quem não tem se apoiava sobretudo em relatos e reportagens (pp.~2--3). Faltava testar de forma
sistemática se a participação na pesquisa de ponta divergiu entre tipos de organização depois de
2012 e se o acesso a computação explica isso. Os autores fazem três perguntas (p.~3): se a pesquisa
ficou mais concentrada, quem são os principais contribuintes e o que aconteceu com quem já atuava
na área.

\subsection*{4. Objetivo do trabalho}
Estimar o efeito causal da ascensão do aprendizado profundo, marcada pelo concurso ImageNet de
2012, sobre a participação de empresas e de universidades de diferentes faixas do \emph{ranking} QS
nas conferências de IA, e mostrar que o acesso a computação é o mecanismo.

\subsection*{5. Hipótese ou pressuposto}
A hipótese é explícita (pp.~4 e 9--10). Partindo de uma função de produção de conhecimento que
depende de conhecimento, habilidades, materiais, equipamento e esforço (eq.~1, p.~9), os autores
preveem que equipamento mais caro reduz a produção de quem não tem acesso a ele: ``we should
observe a divergence between the two groups'' (p.~10), isto é, grandes empresas e universidades de
elite de um lado e as demais universidades do outro. O pressuposto de identificação também é
explícito (pp.~11--12): o sucesso das GPUs no ImageNet 2012 foi inesperado e por isso funciona como
choque exógeno, ligado ao aumento do \emph{compute} disponível e não ao comportamento anterior das
organizações.

\subsection*{6. Base de dados ou corpus}
Artigos de 2000 a 2019 coletados no Scopus, das conferências mais seletivas de computação segundo
o csrankings.org (pp.~12--15). De 63 conferências e 175.491 artigos, ficaram 171.394 artigos de 57
conferências, depois de excluir anos com menos de 25 artigos e a conferência NAACL. As 10
conferências de IA (AAAI, IJCAI, NeurIPS, ICML, KDD, ACL, EMNLP, CVPR, ECCV e ICCV) formam o grupo
tratado e as outras 47, o controle. As afiliações foram classificadas pela Fortune Global 500 de
2018 (46 empresas de tecnologia), pelos \emph{rankings} QS e US News de 2018 e pelo QS de 2011, além
das listas de instituições historicamente negras (HBCU) e hispânicas (HACU). O painel tem 903
observações de conferência por ano (Tabela~1).

\subsection*{7. Metodologia ou abordagem}
O desenho é quase-experimental, com controle sintético generalizado (pp.~10--11). Para cada
conferência de IA, o método usa as conferências de fora da IA para estimar quantos artigos o grupo
teria publicado se 2012 não tivesse acontecido. O efeito (ATT, efeito médio sobre os tratados) é a
diferença média entre o observado e esse contrafactual depois de 2012. O modelo (eq.~2) tem fatores
latentes, efeitos fixos de conferência e de ano e o total de artigos da conferência como controle;
o número de fatores sai de validação cruzada e os erros-padrão, de 2.000 \emph{bootstraps} por
conferência. O desfecho é o número de artigos com ao menos um autor do grupo.

A robustez é testada com placebos de área (Engenharia de Software, Fig.~5) e de data (choque
fictício em 2009, Fig.~6), contagem ponderada pelas afiliações (Tabela~5), \emph{rankings}
alternativos (Tabelas~6 e~7) e um estimador de completamento de matrizes (p.~29). Para o mecanismo,
há uma regressão de efeitos fixos que cruza o choque com uma série anual de \emph{compute} estimada
pela OpenAI (eq.~3, Tabela~8) e uma análise TF-IDF dos resumos, pontuação que mede quanto um termo
é característico dos textos de um grupo (Fig.~8).

\subsection*{8. Principais resultados}
Em relação ao contrafactual, o conjunto das empresas ganhou 46,63 artigos por conferência por ano
(Tabela~2). As 46 grandes empresas de tecnologia ganharam 44,50 (erro-padrão 2,208), contra uma
média de 23 artigos por conferência (Tabela~1). As grandes empresas de fora da tecnologia não
tiveram efeito ($-$0,910), e sem as dez maiores empresas de tecnologia o efeito cai para 19,94. A
Fig.~4 mostra o efeito subindo até cerca de 180 artigos em 2019. Os artigos só de empresas
cresceram cerca de 8 por conferência por ano, e os feitos com universidades, cerca de 42,
concentrados nas universidades do topo do QS (p.~22).

Entre as universidades (Tabela~4), as QS 1--50 ganharam 40,18 (4,49). As faixas 51--100 e
101--200 não tiveram efeito significativo. As 201--300 perderam 14,34 (5,14), e as 301--500
perderam 5,57 (1,54), cerca de 25\% abaixo do contrafactual. As HBCU não tiveram efeito. Os
placebos não mostram efeito (Figs.~5 e~6), e os resultados se mantêm com a contagem ponderada e com
os \emph{rankings} alternativos (Tabelas~5 a~7). Em 20 anos, as HBCU somam 260 artigos e as HACU,
2913; só a Microsoft tem 3302 (pp.~16--17). No mecanismo, a interação entre \emph{compute} e choque é
positiva para empresas e elite e negativa para as universidades QS 101 a 500 (Tabela~8), e as
empresas aparecem mais associadas a termos de aprendizado profundo (Fig.~8).

\subsection*{9. Conclusões dos autores}
A IA está sendo moldada por poucos atores, sobretudo grandes empresas de tecnologia e
universidades de elite, e centenas de universidades de médio e baixo escalão foram deslocadas das
principais conferências (pp.~38--39). Os autores atribuem isso ao acesso desigual a computação, que
chamam de \emph{compute divide}, e dizem ser o primeiro estudo a mostrar que a necessidade de
equipamento especializado pode \emph{desdemocratizar} um campo científico. Defendem intervenção
pública, como uma nuvem nacional de pesquisa e a abertura de dados públicos (p.~38), e afirmam que
o resultado contraria a literatura segundo a qual as empresas vêm recuando da pesquisa pública
(pp.~34--35).

\subsection*{10. Limitações do estudo}
Os autores reconhecem (p.~34) que o estudo cobre só conferências de topo e que o efeito para
empresas nos anos finais pode estar superestimado, se a entrada na IA mudou a participação delas
nas conferências de controle.

Percebi outras. A variável de \emph{compute} da Tabela~8 não é medida por organização: é uma série
anual única, a do maior treinamento conhecido no ano, igual para todos os grupos. A nota~10 (p.~30)
admite que ela ``does not necessarily mean that all organizations need or have access to that
level of computing power'' (p.~30). A escala dessa série também não fecha: o texto diz que ela foi
usada em log (p.~30), mas a Tabela~1 a descreve em escala original, com desvio-padrão de 41.700,14,
e é esse desvio que a p.~31 usa para interpretar os coeficientes.

O controle de total de artigos da conferência (Tabelas~2 e~4) é afetado pelo próprio choque, já
que as conferências de IA cresceram depois de 2012, e isso pode distorcer o efeito. O desfecho é uma
contagem (p.~14), enquanto as Figs.~1 e~2 mostram participações; com conferências de IA maiores,
parte do ganho absoluto das empresas pode ser só crescimento.

Texto e tabelas divergem em vários pontos. A p.~23 fala em 8 artigos a menos para QS 201--300, e a
Tabela~4 mostra $-$14,34. A p.~31 diz que QS 51--100 teve efeito positivo significativo, e a
Tabela~8 mostra $-$0,0003 sem significância. Na mesma p.~31, QS 101--200 aparece primeiro sem
impacto e depois perdendo espaço, e a Tabela~8 pula a coluna~(5). A robustez com completamento de
matrizes fica ``available upon request'' (p.~29), e os outros placebos de área aparecem só na
nota~8 (p.~25). Por fim, o artigo é um \emph{preprint} do arXiv sem revisão por pares na versão de
2020, compara tipos de organização e não países, e as listas HBCU e HACU (p.~16) só valem para os
Estados Unidos.

\subsection*{11. Contribuições}
É a primeira medida sistemática, com desenho contrafactual e placebos, da concentração da pesquisa
em IA depois de 2012. Separa o ganho das grandes empresas de tecnologia do ganho de outras empresas
e das universidades por faixa de \emph{ranking}, e popularizou a expressão \emph{compute divide}.
Para o nosso tema, mostra que o acesso desigual à IA começa antes do produto comercial, já na
produção do conhecimento.

\subsection*{12. Avaliação crítica}
A base é grande e cobre 20 anos (Tabela~1), e o desenho é mais crível que uma comparação simples de
antes e depois: o ajuste antes de 2012 é bom (Fig.~3), há intervalos de confiança (Fig.~4), dois
placebos (Figs.~5 e~6) e três variações de medida e de \emph{ranking} (Tabelas~5 a~7).

\textit{(a) Auditoria da figura principal.} A afirmação central, de que as grandes empresas ganharam
espaço por causa do choque de 2012, está na Fig.~3 (p.~21) e na Tabela~2 (p.~19). Na Fig.~3 o eixo
vertical parte de perto de zero, mas sem o rótulo 0, e não há faixa de incerteza, que só aparece na
Fig.~4. Cada ponto é a média de 10 conferências tratadas, e o ATT médio de 44,50 da Tabela~2 esconde
que quase todo o efeito está nos três últimos anos, quando a série observada vai de cerca de 95
para cerca de 250. O estimador de completamento de matrizes é citado (p.~29), mas não tabulado, e a
participação percentual usada nas Figs.~1 e~2 não aparece na estimação. A Fig.~8 tem outro
problema: não sustenta duas frases da p.~33. No termo \emph{computational cost}, a barra das grandes
empresas é a mais alta e a das QS 301--500, a mais baixa, o contrário do que o texto diz. Em
\emph{support vector machine} e \emph{bayesian}, quem lidera é a elite (QS 1--50), não as
universidades fora dela. Eu esperaria ver o efeito também em participação percentual, com faixa de
confiança, para separar concentração de simples crescimento das conferências de IA.

\textit{(b) Exame de uma afirmação do resumo.} O resumo diz que a divergência entre grandes empresas
e universidades fora da elite ``is driven by access to computing power or compute'' (p.~1). As
únicas peças que tratam disso são a Tabela~8 (p.~30) e a Fig.~8 (p.~34). Na Tabela~8, a série de
\emph{compute} é igual para todos os grupos, então a interação com o choque funciona como uma
tendência pós-2012 das conferências de IA e capta qualquer coisa que tenha crescido no período. A
Fig.~8 mostra sobre o que cada grupo escreve, numa única conferência, e não que \emph{hardware} cada
um tem. Os próprios autores listam três vantagens das empresas que cresceram juntas depois de 2012,
talento, dados proprietários e \emph{compute} (pp.~7--8), e nenhuma análise separa a última das
outras duas. Classificação: ganho de origem não identificada. A divergência está documentada
(Tabelas~2 e~4); a atribuição ao \emph{compute}, e não ao talento ou aos dados, não está.
